% Reviewed translation: PDF pages 61--65 / printed pages 47--51.

\MathBookSourcePage{61}
下面将进一步研究为了唯一确定 $\boldsymbol\phi$ 而必须为这个方程补充的若干边界条件.

\setcounter{statement}{8}
\begin{Remark}\label{rem:sec3-zero-curl-zero-tangential-implies-zero}
Theorem~\ref{thm:sec3-three-dimensional-vector-potential} 的一个直接推论是:
如果 $\mathbf v\in L^2(\Omega)^3$ 满足 \eqref{eq:sec3-divfree-zero-flux},
并且 $\Curl\mathbf v=\mathbf0$ 以及
$\mathbf v\times\mathbf n|_\Gamma=\mathbf0$, 则 $\mathbf v=\mathbf0$.
\end{Remark}

\setcounter{statement}{9}
\begin{Remark}\label{rem:sec3-connected-boundary-divfree-condition}
当 $\Gamma$ 只有一个连通分支(即 $p=0$)时,
\eqref{eq:sec3-divfree-zero-flux} 简化为唯一的条件 $\Div\mathbf v=0$.
\end{Remark}

下面的结果给出关于向量势正则性的更多信息.

\setcounter{statement}{10}
\begin{Remark}\label{rem:sec3-vector-potential-w1s}
当 $\mathbf v\in H\cap L^s(\Omega)^3$ 且 $s>2$ 时, 可以证明
Theorem~\ref{thm:sec3-three-dimensional-vector-potential} 的构造给出
$W^{1,s}(\Omega)^3$ 中的向量势 $\boldsymbol\phi$, 且
$\Div\boldsymbol\phi=0$.
\end{Remark}

\setcounter{statement}{2}
\begin{Corollary}\label{cor:sec3-h1-divfree-h2-potential}
设 $\mathbf v\in H^1(\Omega)^3$ 满足 \eqref{eq:sec3-divfree-zero-flux}.
则它在 $H^2(\Omega)^3$ 中有一个无散向量势 $\boldsymbol\phi$:
\[
\mathbf v=\Curl\boldsymbol\phi,
\qquad
\Div\boldsymbol\phi=0.
\]
\end{Corollary}

证明留作 Exercise.(提示: 使用 Lemma~\ref{lem:sec2-div-free-lifting},
通过匹配 $\mathbf v$ 与其延拓的边界值把 $\mathbf v$ 延拓到 $\mathbb R^3$;
然后利用 Fourier 变换构造向量势.)

\setcounter{statement}{11}
\begin{Remark}\label{rem:sec3-fractional-potential-regularity}
在 Theorem~\ref{thm:sec3-three-dimensional-vector-potential} 与
Corollary~\ref{cor:sec3-h1-divfree-h2-potential} 之间插值, 立即可见:
如果 $\mathbf v\in H^s(\Omega)^3$, $s\in[0,1]$, 满足
\eqref{eq:sec3-divfree-zero-flux}, 则它在 $H^{s+1}(\Omega)^3$ 中有一个无散向量势
$\boldsymbol\phi$.
\end{Remark}

\setcounter{statement}{12}
\begin{Remark}\label{rem:sec3-high-order-potential-regularity}
当 $\Gamma$ 属于 $\mathcal C^{m,1}$ 类($m\ge2$)时, 可以证明每个满足
\eqref{eq:sec3-divfree-zero-flux} 的 $\mathbf v\in H^m(\Omega)^3$
在 $H^{m+1}(\Omega)^3$ 中都有一个无散向量势 $\boldsymbol\phi$.
\end{Remark}

现在转向刻画 $\boldsymbol\phi$ 的边界条件.

\setcounter{statement}{4}
\begin{Theorem}\label{thm:sec3-normal-zero-vector-potential}
$1^\circ)$ 设 $\mathbf v\in L^2(\Omega)^3$ 满足
\eqref{eq:sec3-divfree-zero-flux}. 在满足
\begin{equation}\label{eq:sec3-vector-potential-divfree-system}
\Curl\boldsymbol\phi=\mathbf v,
\qquad
\Div\boldsymbol\phi=0
\end{equation}
的向量势 $\boldsymbol\phi$ 中, 可以选择
$\boldsymbol\phi\in H(\Curl;\Omega)$ 使得
\begin{equation}\label{eq:sec3-vector-potential-normal-zero}
\boldsymbol\phi\cdot\mathbf n=0.
\end{equation}
此外, 当 $\Omega$ 单连通时, $\boldsymbol\phi$ 由
\eqref{eq:sec3-vector-potential-divfree-system} 和
\eqref{eq:sec3-vector-potential-normal-zero} 完全确定; 它可刻画为下列边值问题的唯一解:
\begin{equation}\label{eq:sec3-normal-zero-potential-boundary-problem}
\left\{\begin{aligned}
&\boldsymbol\phi\in H,\\
&-\Delta\boldsymbol\phi=\Curl\mathbf v
&&\text{in }H^{-1}(\Omega)^3,\\
&(\Curl\boldsymbol\phi-\mathbf v)\cdot\mathbf n=0
&&\text{on }\Gamma.
\end{aligned}\right.
\end{equation}
$2^\circ)$ 如果 $\Gamma$ 光滑(例如属于 $\mathcal C^{1,1}$ 类),
则 $\boldsymbol\phi\in H^1(\Omega)^3$.
\end{Theorem}

\MathBookSourcePage{62}
\begin{Proof}
$1^\circ)$ 由 Theorem~\ref{thm:sec3-three-dimensional-vector-potential},
存在满足 \eqref{eq:sec3-vector-potential-divfree-system} 的
$\boldsymbol\phi\in H^1(\Omega)^3$. 因此, 如果能找到
$\boldsymbol\psi\in H(\Curl;\Omega)$ 使得
\begin{equation}\label{eq:sec3-potential-normal-correction}
\left\{\begin{aligned}
\Curl\boldsymbol\psi&=\mathbf0,\\
\Div\boldsymbol\psi&=0
\end{aligned}\right\}\quad\text{in }\Omega,
\qquad
\boldsymbol\psi\cdot\mathbf n=\boldsymbol\phi\cdot\mathbf n
\quad\text{on }\Gamma,
\end{equation}
则 $\boldsymbol\phi-\boldsymbol\psi$ 就是所需函数. 为此, 令
$\chi\in H^1(\Omega)/\mathbb R$ 为下列非齐次 Neumann 问题的解:
\begin{equation}\label{eq:sec3-potential-correction-neumann}
\left\{\begin{aligned}
\Delta\chi&=0&&\text{in }\Gamma,\\
\partial\chi/\partial n&=\boldsymbol\phi\cdot\mathbf n&&\text{on }\Gamma.
\end{aligned}\right.
\end{equation}
则
\[
\boldsymbol\psi=\Grad\chi
\]
恰好满足 \eqref{eq:sec3-potential-normal-correction}, 且当然属于
$H(\Curl;\Omega)$.

当 $\Omega$ 单连通时, $\boldsymbol\phi$ 的唯一性显然是
Remark~\ref{rem:sec2-h-curlzero} 的推论. 此外, 显然 $\boldsymbol\phi$ 满足
\eqref{eq:sec3-normal-zero-potential-boundary-problem}. 反之,
\eqref{eq:sec3-normal-zero-potential-boundary-problem} 至多有一个解, 因为如果其中
$\mathbf v=\mathbf0$, 则一方面 $\Curl\boldsymbol\phi\in H$, 另一方面
$\Curl\boldsymbol\phi\in\operatorname{Ker}(\Curl)$. 因而
$\Curl\boldsymbol\phi=\mathbf0$. 又因 $\boldsymbol\phi\in H$,
Remark~\ref{rem:sec2-h-curlzero} 再次蕴含 $\boldsymbol\phi=\mathbf0$.

$2^\circ)$ 如果 $\Gamma$ 属于 $\mathcal C^{1,1}$ 类, 则
$\boldsymbol\phi\cdot\mathbf n\in H^{1/2}(\Gamma)$. 因此,
Theorem~\ref{thm:laplace-neumann-regularity}(取 $k=0$)断言
$\chi\in H^2(\Omega)$. 从而 $\Grad\chi\in H^1(\Omega)^3$.
\end{Proof}

\setcounter{statement}{5}
\begin{Theorem}\label{thm:sec3-tangential-zero-vector-potential}
$1^\circ)$ 每个满足 \eqref{eq:sec3-divfree-zero-flux} 的
$L^2(\Omega)^3$ 中函数至多有一个无散向量势
$\boldsymbol\phi\in H(\Curl;\Omega)$ 满足边界条件
\begin{equation}\label{eq:sec3-vector-potential-tangential-flux-conditions}
\boldsymbol\phi\times\mathbf n=\mathbf0\quad\text{on }\Gamma,
\qquad
\langle\boldsymbol\phi\cdot\mathbf n,1\rangle_{\Gamma_i}=0,
\quad0\le i\le p.
\end{equation}
$2^\circ)$ 当 $\Omega$ 单连通时, 每个 $\mathbf v\in H$ 恰有一个满足
\eqref{eq:sec3-vector-potential-tangential-flux-conditions} 的无散向量势
$\boldsymbol\phi\in H(\Curl;\Omega)$. 它可刻画为下列边值问题的唯一解:
\begin{equation}\label{eq:sec3-tangential-zero-potential-boundary-problem}
\left\{\begin{aligned}
-\Delta\boldsymbol\phi&=\Curl\mathbf v&&\text{in }H^{-1}(\Omega)^3,\\
\Div\boldsymbol\phi&=0&&\text{in }\Omega,\\
\boldsymbol\phi&\text{ satisfies }\eqref{eq:sec3-vector-potential-tangential-flux-conditions}.
\end{aligned}\right.
\end{equation}
$3^\circ)$ 在上述假设之外, 如果 $\Gamma$ 属于 $\mathcal C^{1,1}$ 类,
或者 $\Omega$ 是凸多面体, 则这个向量势 $\boldsymbol\phi$ 属于 $H^1(\Omega)^3$.
\end{Theorem}
\begin{Proof}
$1^\circ)$ 使用反证法. 假设 $\boldsymbol\phi\in H(\Curl;\Omega)$ 满足
\[
\Curl\boldsymbol\phi=\mathbf0\quad\text{in }\Omega,
\qquad
\boldsymbol\phi\times\mathbf n=\mathbf0\quad\text{on }\Gamma,
\]
\[
\Div\boldsymbol\phi=0\quad\text{in }\Omega,
\qquad
\langle\boldsymbol\phi\cdot\mathbf n,1\rangle_{\Gamma_i}=0
\quad\text{for }0\le i\le p.
\]
则 Remark~\ref{rem:sec3-zero-curl-zero-tangential-implies-zero} 说明
$\boldsymbol\phi=\mathbf0$.

\MathBookSourcePage{63}
$2^\circ)$ 设 $\mathbf v\in H$, 并以 $\widetilde{\mathbf v}$ 表示它在
$\Omega$ 外的零延拓. 则 $\Div\widetilde{\mathbf v}=0$ in $\mathbb R^3$;
由于 $\widetilde{\mathbf v}$ 的支集紧, 可以应用
Theorem~\ref{thm:sec3-three-dimensional-vector-potential}, 即存在
$\boldsymbol\psi\in H^1(\mathbb R^3)^3$ 使得
\[
\widetilde{\mathbf v}=\Curl\boldsymbol\psi,
\qquad
\Div\boldsymbol\psi=0\quad\text{in }\mathbb R^3.
\]
思路是先构造一个散度和旋度均为零、并且与 $\boldsymbol\psi$ 具有相同切向分量的函数
$\mathbf w$. 那么差 $\boldsymbol\psi-\mathbf w$ 将满足
\eqref{eq:sec3-vector-potential-tangential-flux-conditions} 中第一个边界条件.
下一步要证明, 在所有这样的函数 $\mathbf w$ 中, 存在一个在每个 $\Gamma_i$ 上与
$\boldsymbol\psi$ 具有相同平均法向分量的函数; 因而它们的差将满足所有要求.

首先注意到
\[
\Curl\boldsymbol\psi=\mathbf0\quad\text{in }\mathbb R^3-\bar\Omega.
\]
现在取包含 $\bar\Omega$ 的开球 $\mathcal C$, 并如 Figure~\ref{fig:sec3-multiply-connected-domain}
所示, 以 $\Omega_i$ 表示 $\mathcal C-\bar\Omega$ 中由 $\Gamma_i$ 围成的连通分支.
因为 $\Omega$ 单连通, 每个 $\Omega_i$ 也单连通(例如参见 Bernardi [8]), 因此
Theorem~\ref{thm:sec2-curl-free-gradient} 蕴含存在
$q_i\in H^1(\Omega_i)$, 使得
\begin{equation}\label{eq:sec3-exterior-potential-gradients}
\boldsymbol\psi=\Grad q_i\quad\text{in }\Omega_i,
\quad0\le i\le p.
\end{equation}
此外, 因为 $\boldsymbol\psi\in H^1(\mathcal C)^3$,
$q_i\in H^2(\Omega_i)$. 接着令 $\chi\in H^1(\Omega)$ 为非齐次 Dirichlet 问题
\begin{equation}\label{eq:sec3-harmonic-extension-boundary-data}
\left\{\begin{aligned}
\Delta\chi&=0&&\text{in }\Omega,\\
\chi|_{\Gamma_i}&=q_i&&0\le i\le p
\end{aligned}\right.
\end{equation}
的解, 并置
\[
\mathbf q=(q_0,\ldots,q_p).
\]
显然, 映射 $\mathbf q\mapsto\chi$ 属于
$\mathcal L(\prod_{i=0}^pH^2(\Omega_i);H^1(\Omega))$, 且集合
$\{\chi,q_0,q_1,\ldots,q_p\}$ 给出 $\chi$ 在 $H^1(\mathcal C)$ 中的一个延拓.
最后取
\[
\mathbf w=\Grad\chi.
\]
于是
\[
\Div\mathbf w=0,
\qquad
\Curl\mathbf w=\mathbf0\quad\text{in }\Omega,
\]
并且
\[
\mathbf w\times\mathbf n=\Grad\chi\times\mathbf n
=\Grad q_i\times\mathbf n=\boldsymbol\psi\times\mathbf n
\quad\text{on }\Gamma_i,
\quad0\le i\le p.
\]
所以函数 $\boldsymbol\psi-\mathbf w$ 满足
\eqref{eq:sec3-vector-potential-divfree-system} 和
\eqref{eq:sec3-vector-potential-tangential-flux-conditions} 中的第一个条件.

当然, 函数 $\mathbf w$ 不唯一, 因为
\eqref{eq:sec3-exterior-potential-gradients} 只把每个 $q_i$ 确定到一个加法常数.
现在要证明, 可以选择这 $p+1$ 个常数, 使得
\[
\langle(\boldsymbol\psi-\mathbf w)\cdot\mathbf n,1\rangle_{\Gamma_i}=0
\quad\text{for }0\le i\le p.
\]

\MathBookSourcePage{64}
为此, 令 $\mathbf q$ 表示 \eqref{eq:sec3-exterior-potential-gradients} 的解的任意一个
(但固定的)代表元, 并且仅在这里用 $\Delta^{-1}(\mathbf q)$ 表示
\eqref{eq:sec3-harmonic-extension-boundary-data} 的相应解. 对
$\mathbf c\in\mathbb R^{p+1}$, 置
\[
\chi(\mathbf c)=\Delta^{-1}(\mathbf q)+\Delta^{-1}(\mathbf c),
\]
以及
\[
\mathbf w(\mathbf c)=\Grad\chi(\mathbf c)
=\Grad\Delta^{-1}(\mathbf q)+\Grad\Delta^{-1}(\mathbf c).
\]
更仔细地考察 $\Grad\Delta^{-1}(\mathbf c)$ 的性质. 首先注意到, 集合
\[
E=\{\Grad\Delta^{-1}(\mathbf c);\ \mathbf c\in\mathbb R^{p+1}\}
\]
是下列空间的 $p$ 维子空间:
\[
F=\{\mathbf v\in L^2(\Omega)^N;\ \Div\mathbf v=0,\
\ \Curl\mathbf v=\mathbf0,\ \mathbf v\times\mathbf n|_\Gamma=\mathbf0\},
\]
因为 $\Delta^{-1}$ 是同构, 而 $\operatorname{Ker}(\Grad)$ 的维数为 $1$.
接着注意到, 对所有无散 $\mathbf v$ 都有
\[
\langle\mathbf v\cdot\mathbf n,1\rangle_{\Gamma_0}
=-\sum_{i=1}^p\langle\mathbf v\cdot\mathbf n,1\rangle_{\Gamma_i}.
\]
因此, 如果证明映射
\[
\mathbf v\mapsto
(\langle\mathbf v\cdot\mathbf n,1\rangle_{\Gamma_i})_{1\le i\le p}
\]
从 $E$ 到 $\mathbb R^p$ 是单射, 那么因为 $E$ 与 $\mathbb R^p$ 维数相同,
它自动是满射. 但由 $1^\circ)$ 立即可知这个映射确实是单射. 因而可以找到
$\mathbf c\in\mathbb R^{p+1}$, 使得
\[
\langle\mathbf w(\mathbf c)\cdot\mathbf n,1\rangle_{\Gamma_i}
=\langle\boldsymbol\psi\cdot\mathbf n,1\rangle_{\Gamma_i}
\quad\text{for }0\le i\le p,
\]
并且函数 $\boldsymbol\phi=\boldsymbol\psi-\mathbf w(\mathbf c)$
确实满足 \eqref{eq:sec3-vector-potential-divfree-system},
\eqref{eq:sec3-vector-potential-tangential-flux-conditions} 和
\eqref{eq:sec3-tangential-zero-potential-boundary-problem}.

最后证明 \eqref{eq:sec3-tangential-zero-potential-boundary-problem} 至多有一个解.
当 $\mathbf v=\mathbf0$ 时, 由 Remark~\ref{rem:sec2-curl-in-h},
既有 $\Curl\boldsymbol\phi\in\operatorname{Ker}(\Curl)$,
又有 $\Curl\boldsymbol\phi\in H$. 因而 Remark~\ref{rem:sec2-h-curlzero}
蕴含 $\Curl\boldsymbol\phi=\mathbf0$. 所以由
Remark~\ref{rem:sec3-zero-curl-zero-tangential-implies-zero},
$\boldsymbol\phi=\mathbf0$.

$3^\circ)$ 因为 $q_i\in H^2(\Omega_i)$, 由
Theorem~\ref{thm:laplace-dirichlet-regularity}, 当 $\Gamma$ 属于
$\mathcal C^{1,1}$ 类或 $\Omega$ 为凸多面体时,
\eqref{eq:sec3-harmonic-extension-boundary-data} 的解 $\chi$ 属于 $H^2(\Omega)$.
因此 $\mathbf w\in H^1(\Omega)^3$, $\boldsymbol\phi$ 也是如此.
\end{Proof}

\setcounter{statement}{13}
\begin{Remark}\label{rem:sec3-convex-polyhedron-potential-w1p}
设 $\Omega$ 是有界凸多面体, 且 $\mathbf v\in H\cap L^s(\Omega)^3$,
$s>2$. 由 Theorem~\ref{thm:sec3-tangential-zero-vector-potential} 和
Remark~\ref{rem:sec3-vector-potential-w1s}, 其满足
\[
\Div\boldsymbol\phi=0,
\qquad
\boldsymbol\phi\times\mathbf n=\mathbf0\quad\text{on }\Gamma
\]
的唯一向量势 $\boldsymbol\phi$ 还具有正则性
$\boldsymbol\phi\in W^{1,r}(\Omega)^3$, 其中某个 $r$ 满足 $2<r\le s$.
事实上, 通过问题 \eqref{eq:sec3-harmonic-extension-boundary-data},
$\boldsymbol\phi$ 的正则性由右端属于 $L^s(\Omega)$ 的 $-\Delta$ Dirichlet 问题的解的正则性确定
(参见 Theorem~\ref{thm:laplace-dirichlet-regularity}).
\end{Remark}

与二维情形一样, 由 Theorem~\ref{thm:sec3-tangential-zero-vector-potential}
得到 $L^2(\Omega)^3$ 和 $H_0^1(\Omega)^3$ 的如下分解.

\setcounter{statement}{3}
\begin{Corollary}\label{cor:sec3-l2-three-dimensional-decomposition}
$L^2(\Omega)^3$ 中每个函数 $\mathbf v$ 都有正交分解
\begin{equation}
\tag{\ref*{eq:sec3-three-d-fourier-potential-divfree}}
\label{eq:sec3-l2-three-dimensional-decomposition}
\mathbf v=\Grad q+\Curl\boldsymbol\phi,
\end{equation}
\stepcounter{equation}% The source prints 3.22 here but later identifies this object as 3.32.

\MathBookSourcePage{65}
其中 $q\in H^1(\Omega)/\mathbb R$ 是问题
\eqref{eq:sec3-neumann-variational} 的唯一解, 且
$\boldsymbol\phi\in H^1(\Omega)^3$ 满足
$\Curl\boldsymbol\phi\in H$ 和 $\Div\boldsymbol\phi=0$.
此外, 如果 $\Omega$ 单连通, 可以取 $\boldsymbol\phi$ 为
$H(\Curl;\Omega)$ 中问题
\eqref{eq:sec3-tangential-zero-potential-boundary-problem} 的唯一解.
\end{Corollary}

\setcounter{statement}{4}
\begin{Corollary}\label{cor:sec3-h01-three-dimensional-decomposition}
$H_0^1(\Omega)^3$ 中每个函数 $\mathbf v$ 都有正交分解
\begin{equation}\label{eq:sec3-h01-three-dimensional-decomposition}
\mathbf v=(-\Delta)^{-1}\Grad q+\Curl\boldsymbol\phi,
\end{equation}
其中 $\boldsymbol\phi\in H^2(\Omega)^3$, 满足
$\Div\boldsymbol\phi=0$, $\Curl\boldsymbol\phi\in V$, 并且是下列问题的解:
\begin{equation}\label{eq:sec3-h01-three-dimensional-potential-problem}
(-\Delta\boldsymbol\phi,\Curl\mathbf w)
=(\Curl\mathbf v,\Curl\mathbf w)
\quad\forall\mathbf w\in V,
\end{equation}
而 $q\in L_0^2(\Omega)$ 由
\eqref{eq:sec3-h01-three-dimensional-decomposition} 唯一确定.
此外, 当 $\Omega$ 单连通时, 问题
\eqref{eq:sec3-h01-three-dimensional-potential-problem} 在满足
$\boldsymbol\phi\in H$ 和 $\Curl\boldsymbol\phi\in V$ 的函数中恰有一个解.
\end{Corollary}

\setcounter{subsection}{3}
\subsection{$H(\Div;\Omega)\cap H_0(\Curl;\Omega)$ 到 $H^1(\Omega)^3$ 的嵌入}
\label{subsec:sec3-hdiv-h0curl-embedding}

在本节中始终假设 $\Omega$ 有界, 并且 $\Gamma$ 属于 $\mathcal C^{1,1}$ 类,
或者 $\Omega$ 是凸多面体. 与 Subsection~\ref{subsec:sec3-hdiv-hcurl-regularity} 类似, 置
\[
W=H(\Div;\Omega)\cap H_0(\Curl;\Omega),
\qquad
W_0=\{\mathbf w\in W;\ \Div\mathbf w=0\}.
\]

\setcounter{statement}{2}
\begin{Lemma}\label{lem:sec3-w-embedding-reduction}
空间 $W$ 连续嵌入 $H^1(\Omega)^3$, 当且仅当 $W_0$ 也连续嵌入
$H^1(\Omega)^3$.
\end{Lemma}
\begin{Proof}
设 $\boldsymbol\phi\in W$, 并令 $q\in H^1(\Omega)$ 为 Dirichlet 问题
\[
\Delta q=\Div\boldsymbol\phi\quad\text{in }\Omega,
\qquad
q=0\quad\text{on }\Gamma
\]
的解. 则
\[
\|\Grad q-\boldsymbol\phi\|_{0,\Omega}
\le\|\boldsymbol\phi\|_{0,\Omega},
\]
而对 $\Omega$ 的假设蕴含 $q\in H^2(\Omega)$, 且
\[
\|q\|_{2,\Omega}\le C_1\|\Div\boldsymbol\phi\|_{0,\Omega}.
\]
此外, $(\Grad q)\times\mathbf n=\mathbf0$ on $\Gamma$. 因此函数
\[
\boldsymbol\psi=\boldsymbol\phi-\Grad q
\]
属于 $W_0$. 所以, 如果 $W_0$ 连续嵌入 $H^1(\Omega)^3$, 则
$\boldsymbol\phi\in H^1(\Omega)^3$, 并且
\begin{align*}
\|\boldsymbol\phi\|_{1,\Omega}
&\le\|\boldsymbol\psi\|_{1,\Omega}
+C_1\|\Div\boldsymbol\phi\|_{0,\Omega}\\
&\le C_2(\|\boldsymbol\psi\|_{0,\Omega}
+\|\Curl\boldsymbol\psi\|_{0,\Omega})
+C_1\|\Div\boldsymbol\phi\|_{0,\Omega}\\
&\le C_2(\|\boldsymbol\phi\|_{0,\Omega}
+\|\Curl\boldsymbol\phi\|_{0,\Omega})
+C_1\|\Div\boldsymbol\phi\|_{0,\Omega}.
\end{align*}
这就说明该嵌入也是连续的. 反向结论显然.
\end{Proof}
